% v18. G_{5,10}; no denominator cutoff. Depends on mh18:matching.
\section{\texorpdfstring{$G_{5,10}$}{G5,10} 型の完全排除}
\label{gten18:section}

\begin{theorem}[$G_{5,10}$ の障害]
\label{gten18:no-support}
五次の算術的接着データの高分母 support は
\[
 G_{5,10}=\{01,02,03,12,13,24,34\}
\]
と同型ではあり得ない．分母の上限は仮定しない．
証明には正の局所整数データ，marked clutching，非負 nodes，
対向 $D$ 相殺を用いる．
\end{theorem}

\subsection{二価 core と二つの低 row}

頂点 $4$ の高分母近傍 $2,3$ の間は低辺である．従って
補題~\ref{kb15:no-negative-low-wedge} により，core $4$ は正の二高分母枝に属する．
graph の自己同型 $0\leftrightarrow1$，$2\leftrightarrow3$ を使って
canonical な標識を $(A,B,C,D)=(2,3,0,1)$ にそろえる．その行列は
\begin{equation}\label{gten18:central-matrix}
 P^4=
 \begin{pmatrix}
 0&-A&x&-s\\
 -B&0&-y&z\\
 r&-r-1&0&0\\
 -S&T&S-T-1&0
 \end{pmatrix},
\end{equation}
ただし行・列の順は $(2,3,0,1)$，
\[
 \begin{gathered}
 a=q_{24}=sS+1,\quad b=q_{34}=zT+1,\quad
 x=\frac{sS}{r},\quad y=\frac{zT}{r+1},\\
 A=a+x-s,\quad B=b-y+z,\quad AB=ab-1.
 \end{gathered}
\]
全ての $r,s,z,S,T,x,y$ は正整数で，
\[
 S>r,\qquad rT>(r+1)S,\qquad S\ge2.
\]
$\delta=T-S$ と置くと $r\delta>S>r$ なので $\delta\ge2$ である．

低辺 $04,14$ の node 数を
\[
 (n_1,n_2,n_3)=(n_{04;1},n_{04;2},n_{04;3}),\quad
 (m_0,m_2,m_3)=(n_{14;0},n_{14;2},n_{14;3})
\]
と書く．いずれも非負整数で，各組の和は五である．
core $0,1$ は三高分母であり，それらの唯一の低 ray $4$ の rows は
\begin{align}
 R^0&=(p^0_{41},p^0_{42},p^0_{43})
       =(1-n_1, 1-r-n_2, r+2-n_3),\label{gten18:R0}\\
 R^1&=(p^1_{40},p^1_{42},p^1_{43})
       =(\delta+2-m_0, S+1-m_2, 1-T-m_3).
       \label{gten18:R1}
\end{align}
定理~\ref{mh18:matching} により，これらの各 row はちょうど一つの正成分を持つ．
三高分母なので全成分は非零であり，特に
$R^0_2<0$，$R^1_3<0$ である．

$q=q_{01}>1$ とする．共通低 puncture $(01;4)$ の signed-divisor 分類により，
$(R^0_1,R^1_0)$ は通常型 $(u,-u)$，正例外型 $\{1,q-1\}$，
負例外型 $\{-1,-(q-1)\}$ のいずれかである．
以下で三場合を排除する．

\subsection{通常接着の場合}

まず $R^0_1>0$ とする．式~\eqref{gten18:R0} から $R^0_1=1,n_1=0$ であり，
$R^1_0=-1$ から $m_0=\delta+3\le5$ を得る．従って
$\delta=2,m_0=5,m_2=m_3=0$ である．
matching 性により $R^0_3<0$ なので $n_3\ge r+3$，従って $r\le2$ である．
$r<S<2r$ より，唯一の可能性は
\begin{equation}\label{gten18:ordinary-central}
 (r,S,T)=(2,3,5)
\end{equation}
である．このとき $R^0=(1,-1,-1)$，$R^1=(-1,4,-4)$ となる．

次に $R^0_1<0$ とすると $R^1_0>0$ である．matching 性から
$S+1-m_2<0$，すなわち $m_2\ge S+2$ なので $S\le3$ である．
また $R^1_0=-R^0_1\le4$ と $m_0\le3-S$ より
\[
 \delta\le m_0+2\le5-S.
\]
$S=2$ なら $r=1$，正の枝から $\delta\ge3$ であるため
$\delta=3,m_0=1,m_2=4,m_3=0,R^1_0=4$ となる．
すると $n_1=5,n_2=n_3=0$ なので $R^0_2=0$ となり不可能である．
$S=3$ なら $\delta=2,m_0=m_3=0,m_2=5$ であり，
再び式~\eqref{gten18:ordinary-central} を得る．
この場合の rows は $R^0=(-4,-1,4)$，$R^1=(4,-1,-4)$ である．

式~\eqref{gten18:ordinary-central} の整数条件から $s=2u,z=3v$ と書ける．
pair equation は
\[
 (7u+1)(13v+1)=(6u+1)(15v+1)-1,
 \qquad (u-2)(v+1)=-3
\]
となる．$u,v\ge1$ より $u=1,v=2$，従って
\[
 (r,s,z,S,T)=(2,2,6,3,5),\qquad (a,b)=(7,31),\qquad B=27.
\]
しかし $q_{23}=1$ なので，core $3$ の $p^3_{42}$ は $30$ の非零 signed divisor である．
辺 $34$ の puncture $2$ の clutching は
$p^3_{42}-27\equiv0\pmod{31}$ を要求する．
$|p^3_{42}|\le30$ の下でその候補は $27,-4$ だけで，どちらも $30$ を割らない．
これは矛盾である．この通常型の排除には $D$ 相殺を用いていない．

\subsection{正の例外接着の場合}

$R^0_1\le1$ なので，正例外の向きは
\[
 R^0_1=1,\quad R^1_0=q-1,\quad n_1=0,\quad
 q=\delta+3-m_0
\]
である．matching 性と row sum により $R^0=(1,-1,-1)$ となる．
行列~\eqref{gten18:central-matrix} の低 rays における係数は
\begin{equation}\label{gten18:D-low}
 D^4_0=\delta+y-x,\qquad D^4_1=s-z-1.
\end{equation}
辺 $04$ の $D$ 相殺を，core $0$ の低 column を用いて書くと
\[
 q_{02}+q_{03}=5-m_0+x-y.
\]
左辺は六以上なので $x-y>0$ である．
辺 $14$ の $D$ 相殺は
\[
 p^1_{04}+p^1_{24}+p^1_{34}=s-z-2
\]
を与える．ここで $p^1_{04}=1$，他の二成分は負整数なので $s-z\le1$ である．
一方，正の枝の不等式から
\[
 x-y=\frac{sS}{r}-\frac{zT}{r+1}
       <\frac{(s-z)S}{r}
\]
であり，$x-y>0$ は $s-z\ge1$ を要求する．従って $s=z+1$ である．
上の低 column の和は $-1$ となり，二つの負成分はともに $-1$ でなければならない．
pair equations と低 row の和から
\[
 q_{12}+q_{13}=2-(p^1_{42}+p^1_{43})=q+2
\]
を得るが，左辺は偶数，右辺は奇数である．これが矛盾である．

\subsection{負の例外接着と差二の共鳴}

負例外の場合，式~\eqref{gten18:R1} と $\delta\ge2$ から
$R^1_0\ge-1$ であるため，必ず
\begin{equation}\label{gten18:negative-exception}
 R^1_0=-1,\quad R^0_1=-(q-1),\quad
 \delta=2,\quad m_0=5,\quad m_2=m_3=0,\quad
 n_1=q,\quad q\in\{3,5\}
\end{equation}
となる．$h=S-r$ と置けば $1\le h\le r-1$，特に $r\ge2$ である．
正整数
\[
 \alpha=x-s=\frac{sh}{r},\qquad
 \beta=y-z=\frac{z(h+1)}{r+1}
\]
を導入する．二高分母の pair equation は
\begin{equation}\label{gten18:E}
 E\alpha\beta=\alpha-\beta+1,\qquad
 E=\frac{r^2}{h}-\frac{(r+1)^2}{h+1}
\end{equation}
を与える．

辺 $24$ の puncture $3$ では，$p^2_{43}$ は $a-1$ の signed divisor で，
$p^2_{43}\equiv\alpha\pmod a$ である．
$0<\alpha<(a-1)/2$ なので候補は $\alpha$ と $-(a-\alpha)$ だけである．
後者の絶対値は $(a-1)/2$ より大きく，$a-1$ を割るなら $a-1$ 自身である．
従って後者は $\alpha=1$ の場合に限られる．
同様に，辺 $34$ の puncture $2$ では，$p^3_{42}$ は $b-1$ の signed divisor で，
通常の値 $-\beta$ 以外が可能なのは $\beta=1$ のときだけである．

$\alpha,\beta>1$ なら，これらの整除性から
$r^2/h$，$(r+1)^2/(h+1)$ はともに整数となる．
式~\eqref{gten18:E} の右辺を $\alpha\beta$ で割った値は絶対値一未満なので，
$E=0$ である．これは $h=r^2/(2r+1)$ を要求するが，
$\gcd(r,2r+1)=1$ より不可能である．
$\alpha=1,\beta>1$ の場合も $r=sh$ により $r^2/h$ は整数であり，
他方の通常整除性と合わせて $E$ は整数となる．
$E=2/\beta-1$ なので $\beta=2,E=0$ となり，やはり不可能である．
従って必ず $\beta=1$ である．

このとき $E=1$，$r+1=z(h+1)$，$z\ge2$ である．これらを代入すると
\[
 (h+1)z(z-2)=h-1
\]
を得る．$z\ge3$ は不可能であり，$z=2,h=1$ に限られる．
以上により全ての残存模型は，探索上限なしに
\begin{equation}\label{gten18:resonance}
 r=3,\quad S=4,\quad T=6,\quad z=2,\quad s=3\alpha,\quad
 a=12\alpha+1,\quad b=13,\qquad \alpha\in\ZZ_{>0}
\end{equation}
へ帰着した．これは固定分母十三の既知の共鳴族と一致する．

\subsection{共鳴族の \texorpdfstring{$D$}{D} と node による排除}

式~\eqref{gten18:negative-exception}--\eqref{gten18:resonance} により
\[
 R^1=(-1,5,-5).
\]
その低 column を $(-(q-1),U,-V)$ と書けば
$q_{12}=5U+1$，$q_{13}=5V+1$ であり，$U,V$ は正の偶整数である．
式~\eqref{gten18:D-low} と辺 $14$ の $D$ 相殺は
\begin{equation}\label{gten18:UV}
 U-V=3\alpha+q-5
\end{equation}
を与える．従って $\alpha$ は偶数で，$\alpha\ge2$，
$U\ge3\alpha$，$q_{12}\ge15\alpha+1$ である．

core $0$ の低 row は
$R^0=(-(q-1),-(2+n_2),5-n_3)$ である．
初めの二成分は負なので最後は正であり，その低 column は
$(-1,-u,v)$，$u,v>0$ と書ける．辺 $04$ の $D$ 相殺から
\begin{equation}\label{gten18:uv}
 u=v+4\alpha-5\ge4\alpha-4,\qquad
 q_{02}=(2+n_2)u+1\ge8\alpha-7
\end{equation}
を得る．

$\alpha\ge2$ なので，上の signed-divisor 分類は
$p^2_{43}=\alpha$ を強制する．従って $p^2_{34}=12$ である．
core $2$ は三高分母であり，matching 性によって低 ray $3$ は高 ray $4$ と対になる．
そのため $p^2_{40},p^2_{41}$ はともに負である．
辺 $24$ の puncture $3$ の node は二なので，残る二つの node を $\nu_0,\nu_1$ とすると
$\nu_0+\nu_1=3$ である．clutching と node の式は
\[
 p^2_{40}=(8-12\nu_0)\alpha+1-\nu_0,\qquad
 p^2_{41}=(12\nu_0-21)\alpha+\nu_0-2
\]
を与える．両方が負になるのは $\nu_0=1,\nu_1=2$ の場合だけであり，
\[
 p^2_{40}=-4\alpha,\qquad p^2_{41}=-(9\alpha+1)
\]
である．正整数
\[
 M=-p^2_{04}=\frac{a q_{02}-1}{4\alpha}>3q_{02},\qquad
 N=-p^2_{14}=\frac{a q_{12}-1}{9\alpha+1}>q_{12}
\]
を考える．core $4$ の高 ray $2$ の係数は
$D^4_2=12-12\alpha$ なので，辺 $24$ の $D$ 相殺は
\[
 M+N=24\alpha+1
\]
を要求する．一方，式~\eqref{gten18:UV},\eqref{gten18:uv} から
\[
 M+N>3q_{02}+q_{12}
       \ge39\alpha-20
       \ge24\alpha+10
\]
となり矛盾する．これで三つの接着型を全て排除し，
定理~\ref{gten18:no-support} が従う．

この証明で有限に数えたものは node の三成分配分および符号の場合分けである．
分母を切る探索ではなく，残り得る無限族~\eqref{gten18:resonance} も不等式で排除した．
$D$ 相殺を実際に用いた辺は，canonical 標識で $04,14,24$ の三本だけである．
